\nonstopmode{}
\documentclass[a4paper]{book}
\usepackage[times,inconsolata,hyper]{Rd}
\usepackage{makeidx}
\makeatletter\@ifl@t@r\fmtversion{2018/04/01}{}{\usepackage[utf8]{inputenc}}\makeatother
% \usepackage{graphicx} % @USE GRAPHICX@
\makeindex{}
\begin{document}
\chapter*{}
\begin{center}
{\textbf{\huge Package `seroCOP'}}
\par\bigskip{\large \today}
\end{center}
\ifthenelse{\boolean{Rd@use@hyper}}{\hypersetup{pdftitle = {seroCOP: Correlates of Protection Analysis Using Bayesian Methods}}}{}
\ifthenelse{\boolean{Rd@use@hyper}}{\hypersetup{pdfauthor = {David Hodgson}}}{}
\begin{description}
\raggedright{}
\item[Title]\AsIs{Correlates of Protection Analysis Using Bayesian Methods}
\item[Version]\AsIs{0.1.0}
\item[Description]\AsIs{Fits generalized logistic functions to serological data (antibody titres)
and infection outcomes using Bayesian methods implemented via brms. Provides
tools for assessing correlates of risk/protection including ROC AUC, LOO-CV,
and Brier scores. Built around an R6 class for streamlined workflow.}
\item[License]\AsIs{MIT + file LICENSE}
\item[URL]\AsIs{}\url{https://seroanalytics.org/seroCOP/}\AsIs{}
\item[BugReports]\AsIs{}\url{https://github.com/seroanalytics/seroCOP/issues}\AsIs{}
\item[Encoding]\AsIs{UTF-8}
\item[Roxygen]\AsIs{list(markdown = TRUE)}
\item[RoxygenNote]\AsIs{7.3.2}
\item[Depends]\AsIs{R (>= 3.5.0)}
\item[Imports]\AsIs{brms (>= 2.19.0), rstan (>= 2.21.0), rstantools (>= 2.1.1), R6
(>= 2.5.0), pROC (>= 1.18.0), loo (>= 2.4.0), ggplot2 (>=
3.3.0), tidyr (>= 1.0.0), methods, stats}
\item[LinkingTo]\AsIs{BH (>= 1.66.0), Rcpp (>= 0.12.0), RcppEigen (>= 0.3.3.3.0),
RcppParallel (>= 5.0.1), rstan (>= 2.21.0), StanHeaders (>=
2.21.0)}
\item[Suggests]\AsIs{knitr, rmarkdown, pkgdown, testthat (>= 3.0.0)}
\item[VignetteBuilder]\AsIs{knitr}
\item[Config/testthat/edition]\AsIs{3}
\item[Biarch]\AsIs{true}
\item[SystemRequirements]\AsIs{GNU make, C++14}
\item[NeedsCompilation]\AsIs{no}
\item[Author]\AsIs{David Hodgson [aut, cre]}
\item[Maintainer]\AsIs{David Hodgson }\email{david.hodgson@example.com}\AsIs{}
\end{description}
\Rdcontents{Contents}
\HeaderA{seroCOP-package}{seroCOP: Correlates of Protection Analysis Using Bayesian Methods}{seroCOP.Rdash.package}
\aliasA{seroCOP}{seroCOP-package}{seroCOP}
\keyword{internal}{seroCOP-package}
%
\begin{Description}
Fits generalized logistic functions to serological data (antibody titres) and infection outcomes using Bayesian methods implemented via brms. Provides tools for assessing correlates of risk/protection including ROC AUC, LOO-CV, and Brier scores. Built around an R6 class for streamlined workflow.
\end{Description}
%
\begin{Author}
\strong{Maintainer}: David Hodgson \email{david.hodgson@example.com}

\end{Author}
%
\begin{SeeAlso}
Useful links:
\begin{itemize}

\item{} \url{https://seroanalytics.org/seroCOP/}
\item{} Report bugs at \url{https://github.com/seroanalytics/seroCOP/issues}

\end{itemize}


\end{SeeAlso}
\HeaderA{brier\_score}{Calculate Brier score}{brier.Rul.score}
%
\begin{Description}
Calculate the Brier score for probabilistic predictions
\end{Description}
%
\begin{Usage}
\begin{verbatim}
brier_score(observed, predicted)
\end{verbatim}
\end{Usage}
%
\begin{Arguments}
\begin{ldescription}
\item[\code{observed}] Binary vector of observed outcomes (0/1)

\item[\code{predicted}] Vector of predicted probabilities
\end{ldescription}
\end{Arguments}
%
\begin{Value}
Numeric Brier score (lower is better)
\end{Value}
%
\begin{Examples}
\begin{ExampleCode}
observed <- c(0, 1, 1, 0, 1)
predicted <- c(0.1, 0.8, 0.7, 0.2, 0.9)
brier_score(observed, predicted)
\end{ExampleCode}
\end{Examples}
\HeaderA{extract\_parameters}{Extract parameter estimates from SeroCOP fit}{extract.Rul.parameters}
%
\begin{Description}
Extract posterior summaries of model parameters
\end{Description}
%
\begin{Usage}
\begin{verbatim}
extract_parameters(serocop_obj, prob = 0.95)
\end{verbatim}
\end{Usage}
%
\begin{Arguments}
\begin{ldescription}
\item[\code{serocop\_obj}] A fitted SeroCOP object

\item[\code{prob}] Probability for credible intervals (default 0.95)
\end{ldescription}
\end{Arguments}
%
\begin{Value}
Data frame with parameter estimates and credible intervals
\end{Value}
\HeaderA{SeroCOP}{SeroCOP R6 Class}{SeroCOP}
\keyword{r6-classes}{SeroCOP}
%
\begin{Description}
An R6 class for analyzing correlates of protection using Bayesian methods.
Fits a four-parameter logistic model to antibody titre and infection outcome data.

This class implements a Bayesian logistic model for analyzing correlates of protection.
It provides methods for model fitting, prediction, performance evaluation, and visualization.
\end{Description}
%
\begin{Details}
SeroCOP R6 Class for Correlates of Protection Analysis

This class provides a complete workflow for correlates of protection analysis:
\begin{enumerate}

\item{} Data input and validation
\item{} Bayesian model fitting using brms (which uses Stan backend)
\item{} Model diagnostics and validation
\item{} Performance metrics (ROC AUC, Brier score, LOO-CV)
\item{} Visualization

\end{enumerate}

\end{Details}
%
\begin{Section}{Public fields}

\begin{description}

\item[\code{titre}] Numeric vector of antibody titres (log scale recommended)

\item[\code{infected}] Binary vector of infection status (0/1)

\item[\code{group}] Optional factor vector for hierarchical modeling

\item[\code{fit}] Stan fit object (after fitting)

\item[\code{loo}] LOO-CV object (after fitting)

\item[\code{priors}] List of prior distributions for model parameters

\item[\code{weights}] Optional numeric vector of observation weights

\end{description}


\end{Section}
%
\begin{Section}{Methods}
%
\begin{SubSection}{Public methods}
\begin{itemize}

\item{} \Rhref{\#method-SeroCOP-new}{\code{SeroCOP\$new()}}
\item{} \Rhref{\#method-SeroCOP-definePrior}{\code{SeroCOP\$definePrior()}}
\item{} \Rhref{\#method-SeroCOP-fit_model}{\code{SeroCOP\$fit\_model()}}
\item{} \Rhref{\#method-SeroCOP-predict}{\code{SeroCOP\$predict()}}
\item{} \Rhref{\#method-SeroCOP-predict_protection}{\code{SeroCOP\$predict\_protection()}}
\item{} \Rhref{\#method-SeroCOP-extract_group_parameters}{\code{SeroCOP\$extract\_group\_parameters()}}
\item{} \Rhref{\#method-SeroCOP-summary}{\code{SeroCOP\$summary()}}
\item{} \Rhref{\#method-SeroCOP-get_metrics}{\code{SeroCOP\$get\_metrics()}}
\item{} \Rhref{\#method-SeroCOP-plot_curve}{\code{SeroCOP\$plot\_curve()}}
\item{} \Rhref{\#method-SeroCOP-plot_roc}{\code{SeroCOP\$plot\_roc()}}
\item{} \Rhref{\#method-SeroCOP-plot_group_curves}{\code{SeroCOP\$plot\_group\_curves()}}
\item{} \Rhref{\#method-SeroCOP-plot_posteriors}{\code{SeroCOP\$plot\_posteriors()}}
\item{} \Rhref{\#method-SeroCOP-extract_cop}{\code{SeroCOP\$extract\_cop()}}
\item{} \Rhref{\#method-SeroCOP-clone}{\code{SeroCOP\$clone()}}

\end{itemize}

\end{SubSection}



\hypertarget{method-SeroCOP-new}{}
%
\begin{SubSection}{Method \code{new()}}
Create a new SeroCOP object
%
\begin{SubSubSection}{Usage}
\begin{alltt}SeroCOP$new(titre, infected, group = NULL, weights = NULL)\end{alltt}

\end{SubSubSection}


%
\begin{SubSubSection}{Arguments}

\begin{description}

\item[\code{titre}] Numeric vector of antibody titres

\item[\code{infected}] Binary vector (0/1) of infection outcomes

\item[\code{group}] Optional factor or character vector for grouping (e.g., age groups)

\item[\code{weights}] Optional numeric vector of non-negative observation weights.
When supplied, a weighted Bernoulli likelihood is used.
If \code{NULL} (the default) all observations receive equal weight.

\end{description}


\end{SubSubSection}

%
\begin{SubSubSection}{Returns}
A new SeroCOP object
\end{SubSubSection}

\end{SubSection}



\hypertarget{method-SeroCOP-definePrior}{}
%
\begin{SubSection}{Method \code{definePrior()}}
Define custom prior distributions for model parameters
%
\begin{SubSubSection}{Usage}
\begin{alltt}SeroCOP$definePrior(
  floor_alpha = NULL,
  floor_beta = NULL,
  ceiling_alpha = NULL,
  ceiling_beta = NULL,
  ec50_mean = NULL,
  ec50_sd = NULL,
  slope_mean = NULL,
  slope_sd = NULL
)\end{alltt}

\end{SubSubSection}


%
\begin{SubSubSection}{Arguments}

\begin{description}

\item[\code{floor\_alpha}] Alpha parameter for floor beta prior (default: 1)

\item[\code{floor\_beta}] Beta parameter for floor beta prior (default: 9)

\item[\code{ceiling\_alpha}] Alpha parameter for ceiling beta prior (default: 9)

\item[\code{ceiling\_beta}] Beta parameter for ceiling beta prior (default: 1)

\item[\code{ec50\_mean}] Mean for ec50 normal prior (default: midpoint of titre range)

\item[\code{ec50\_sd}] SD for ec50 normal prior (default: titre range / 4)

\item[\code{slope\_mean}] Mean for slope normal prior (default: 0)

\item[\code{slope\_sd}] SD for slope normal prior (default: 2)

\end{description}


\end{SubSubSection}

%
\begin{SubSubSection}{Returns}
Self (invisibly)
\end{SubSubSection}

%
\begin{SubSubSection}{Examples}

\begin{alltt}\bsl{}dontrun\{
model <- SeroCOP$new(titre = titre, infected = infected)

# Use default priors centered on data
model$definePrior()

# Custom priors
model$definePrior(
  floor_alpha = 2, floor_beta = 18,    # Stronger prior for low floor
  ceiling_alpha = 18, ceiling_beta = 2, # Stronger prior for high ceiling
  ec50_mean = 2.0, ec50_sd = 1.0,      # Custom ec50 prior
  slope_mean = 0, slope_sd = 1         # More conservative slope
)
\}
\end{alltt}



\end{SubSubSection}


\end{SubSection}



\hypertarget{method-SeroCOP-fit_model}{}
%
\begin{SubSection}{Method \code{fit\_model()}}
Fit the Bayesian logistic model
%
\begin{SubSubSection}{Usage}
\begin{alltt}SeroCOP$fit_model(chains = 4, iter = 2000, warmup = 1000, thin = 1, ...)\end{alltt}

\end{SubSubSection}


%
\begin{SubSubSection}{Arguments}

\begin{description}

\item[\code{chains}] Number of MCMC chains (default: 4)

\item[\code{iter}] Number of iterations per chain (default: 2000)

\item[\code{warmup}] Number of warmup iterations (default: iter/2)

\item[\code{thin}] Thinning interval (default: 1)

\item[\code{...}] Additional arguments passed to brms::brm

\end{description}


\end{SubSubSection}

%
\begin{SubSubSection}{Returns}
Self (invisibly)
\end{SubSubSection}

\end{SubSection}



\hypertarget{method-SeroCOP-predict}{}
%
\begin{SubSection}{Method \code{predict()}}
Get posterior predictions for infection probability
%
\begin{SubSubSection}{Usage}
\begin{alltt}SeroCOP$predict(newdata = NULL)\end{alltt}

\end{SubSubSection}


%
\begin{SubSubSection}{Arguments}

\begin{description}

\item[\code{newdata}] Optional new titre values for prediction

\end{description}


\end{SubSubSection}

%
\begin{SubSubSection}{Returns}
Matrix of posterior predictions (iterations x observations)
\end{SubSubSection}

\end{SubSection}



\hypertarget{method-SeroCOP-predict_protection}{}
%
\begin{SubSection}{Method \code{predict\_protection()}}
Extract probability of protection from the fitted model
%
\begin{SubSubSection}{Usage}
\begin{alltt}SeroCOP$predict_protection(newdata = NULL)\end{alltt}

\end{SubSubSection}


%
\begin{SubSubSection}{Arguments}

\begin{description}

\item[\code{newdata}] Optional vector of new titre values for prediction

\end{description}


\end{SubSubSection}

%
\begin{SubSubSection}{Returns}
Matrix of protection probabilities (rows = MCMC samples, cols = observations)
\end{SubSubSection}

\end{SubSection}



\hypertarget{method-SeroCOP-extract_group_parameters}{}
%
\begin{SubSection}{Method \code{extract\_group\_parameters()}}
Extract group-specific parameters from hierarchical model
%
\begin{SubSubSection}{Usage}
\begin{alltt}SeroCOP$extract_group_parameters()\end{alltt}

\end{SubSubSection}


%
\begin{SubSubSection}{Returns}
A data.frame with group-specific parameter estimates (NULL if not hierarchical)
\end{SubSubSection}

\end{SubSection}



\hypertarget{method-SeroCOP-summary}{}
%
\begin{SubSection}{Method \code{summary()}}
Get summary statistics for model parameters
%
\begin{SubSubSection}{Usage}
\begin{alltt}SeroCOP$summary()\end{alltt}

\end{SubSubSection}


%
\begin{SubSubSection}{Returns}
Data frame with parameter summaries
\end{SubSubSection}

\end{SubSection}



\hypertarget{method-SeroCOP-get_metrics}{}
%
\begin{SubSection}{Method \code{get\_metrics()}}
Calculate performance metrics
%
\begin{SubSubSection}{Usage}
\begin{alltt}SeroCOP$get_metrics()\end{alltt}

\end{SubSubSection}


%
\begin{SubSubSection}{Returns}
List containing ROC AUC, Brier score, and LOO-CV metrics
\end{SubSubSection}

\end{SubSection}



\hypertarget{method-SeroCOP-plot_curve}{}
%
\begin{SubSection}{Method \code{plot\_curve()}}
Plot the fitted curve with uncertainty
%
\begin{SubSubSection}{Usage}
\begin{alltt}SeroCOP$plot_curve(title = "Correlates of Risk Curve", ...)\end{alltt}

\end{SubSubSection}


%
\begin{SubSubSection}{Arguments}

\begin{description}

\item[\code{title}] Plot title

\item[\code{...}] Additional arguments passed to ggplot2

\end{description}


\end{SubSubSection}

%
\begin{SubSubSection}{Returns}
A ggplot object
\end{SubSubSection}

\end{SubSection}



\hypertarget{method-SeroCOP-plot_roc}{}
%
\begin{SubSection}{Method \code{plot\_roc()}}
Plot ROC curve
%
\begin{SubSubSection}{Usage}
\begin{alltt}SeroCOP$plot_roc(title = "ROC Curve")\end{alltt}

\end{SubSubSection}


%
\begin{SubSubSection}{Arguments}

\begin{description}

\item[\code{title}] Plot title

\end{description}


\end{SubSubSection}

%
\begin{SubSubSection}{Returns}
A ggplot object
\end{SubSubSection}

\end{SubSection}



\hypertarget{method-SeroCOP-plot_group_curves}{}
%
\begin{SubSection}{Method \code{plot\_group\_curves()}}
Plot group-specific curves for hierarchical models
%
\begin{SubSubSection}{Usage}
\begin{alltt}SeroCOP$plot_group_curves(title = "Group-Specific Correlates of Risk")\end{alltt}

\end{SubSubSection}


%
\begin{SubSubSection}{Arguments}

\begin{description}

\item[\code{title}] Plot title

\end{description}


\end{SubSubSection}

%
\begin{SubSubSection}{Returns}
A ggplot object (returns NULL if not a hierarchical model)
\end{SubSubSection}

\end{SubSection}



\hypertarget{method-SeroCOP-plot_posteriors}{}
%
\begin{SubSection}{Method \code{plot\_posteriors()}}
Plot posterior distributions of parameters
%
\begin{SubSubSection}{Usage}
\begin{alltt}SeroCOP$plot_posteriors()\end{alltt}

\end{SubSubSection}


%
\begin{SubSubSection}{Returns}
A ggplot object
\end{SubSubSection}

\end{SubSection}



\hypertarget{method-SeroCOP-extract_cop}{}
%
\begin{SubSection}{Method \code{extract\_cop()}}
Extract the correlate of protection conditional on exposure.
%
\begin{SubSubSection}{Usage}
\begin{alltt}SeroCOP$extract_cop(correlate_of_risk, upper_bound = 0.7)\end{alltt}

\end{SubSubSection}


%
\begin{SubSubSection}{Arguments}

\begin{description}

\item[\code{correlate\_of\_risk}] Numeric vector of correlates of risk.

\item[\code{upper\_bound}] Numeric value for the upper bound (default: 0.7).

\end{description}


\end{SubSubSection}

%
\begin{SubSubSection}{Returns}
Numeric vector of correlates of protection.
\end{SubSubSection}

\end{SubSection}



\hypertarget{method-SeroCOP-clone}{}
%
\begin{SubSection}{Method \code{clone()}}
The objects of this class are cloneable with this method.
%
\begin{SubSubSection}{Usage}
\begin{alltt}SeroCOP$clone(deep = FALSE)\end{alltt}

\end{SubSubSection}


%
\begin{SubSubSection}{Arguments}

\begin{description}

\item[\code{deep}] Whether to make a deep clone.

\end{description}


\end{SubSubSection}

\end{SubSection}

\end{Section}
%
\begin{Examples}
\begin{ExampleCode}
## Not run: 
# Create synthetic data
set.seed(123)
n <- 200
titre <- rnorm(n, mean = 2, sd = 1.5)
prob <- 0.05 + 0.9 / (1 + exp(2 * (titre - 1.5)))
infected <- rbinom(n, 1, prob)

# Initialize and fit
model <- SeroCOP$new(titre = titre, infected = infected)
model$fit(chains = 4, iter = 2000)

# Get metrics
model$get_metrics()

# Plot results
model$plot_curve()
model$plot_roc()

## End(Not run)

## ------------------------------------------------
## Method `SeroCOP$definePrior`
## ------------------------------------------------

## Not run: 
model <- SeroCOP$new(titre = titre, infected = infected)

# Use default priors centered on data
model$definePrior()

# Custom priors
model$definePrior(
  floor_alpha = 2, floor_beta = 18,    # Stronger prior for low floor
  ceiling_alpha = 18, ceiling_beta = 2, # Stronger prior for high ceiling
  ec50_mean = 2.0, ec50_sd = 1.0,      # Custom ec50 prior
  slope_mean = 0, slope_sd = 1         # More conservative slope
)

## End(Not run)
\end{ExampleCode}
\end{Examples}
\HeaderA{SeroCOPMulti}{SeroCOPMulti R6 Class for Multi-Biomarker Correlates of Protection Analysis}{SeroCOPMulti}
\keyword{r6-classes}{SeroCOPMulti}
%
\begin{Description}
An R6 class for analyzing multiple biomarkers simultaneously.
Fits separate models for each biomarker using brms and provides comparison tools.
Supports hierarchical modeling with group-level effects.
\end{Description}
%
\begin{Section}{Public fields}

\begin{description}

\item[\code{titre}] Matrix of antibody titres (columns = biomarkers)

\item[\code{infected}] Binary vector of infection status (0/1)

\item[\code{biomarker\_names}] Names of biomarkers

\item[\code{models}] List of fitted SeroCOP objects

\item[\code{group}] Optional factor vector for hierarchical modeling

\item[\code{weights}] Optional numeric matrix of observation weights (rows = observations, cols = biomarkers)

\end{description}


\end{Section}
%
\begin{Section}{Methods}
%
\begin{SubSection}{Public methods}
\begin{itemize}

\item{} \Rhref{\#method-SeroCOPMulti-new}{\code{SeroCOPMulti\$new()}}
\item{} \Rhref{\#method-SeroCOPMulti-fit_all}{\code{SeroCOPMulti\$fit\_all()}}
\item{} \Rhref{\#method-SeroCOPMulti-compare_biomarkers}{\code{SeroCOPMulti\$compare\_biomarkers()}}
\item{} \Rhref{\#method-SeroCOPMulti-plot_comparison}{\code{SeroCOPMulti\$plot\_comparison()}}
\item{} \Rhref{\#method-SeroCOPMulti-plot_all_curves}{\code{SeroCOPMulti\$plot\_all\_curves()}}
\item{} \Rhref{\#method-SeroCOPMulti-extract_group_parameters}{\code{SeroCOPMulti\$extract\_group\_parameters()}}
\item{} \Rhref{\#method-SeroCOPMulti-plot_group_curves}{\code{SeroCOPMulti\$plot\_group\_curves()}}
\item{} \Rhref{\#method-SeroCOPMulti-extract_cop_multi}{\code{SeroCOPMulti\$extract\_cop\_multi()}}
\item{} \Rhref{\#method-SeroCOPMulti-clone}{\code{SeroCOPMulti\$clone()}}

\end{itemize}

\end{SubSection}



\hypertarget{method-SeroCOPMulti-new}{}
%
\begin{SubSection}{Method \code{new()}}
Create a new SeroCOPMulti object
%
\begin{SubSubSection}{Usage}
\begin{alltt}SeroCOPMulti$new(
  titre,
  infected,
  biomarker_names = NULL,
  group = NULL,
  weights = NULL
)\end{alltt}

\end{SubSubSection}


%
\begin{SubSubSection}{Arguments}

\begin{description}

\item[\code{titre}] Matrix of antibody titres (rows = samples, cols = biomarkers)

\item[\code{infected}] Binary vector (0/1) of infection outcomes

\item[\code{biomarker\_names}] Optional vector of biomarker names

\item[\code{group}] Optional grouping variable for hierarchical modeling

\item[\code{weights}] Optional numeric matrix of non-negative observation weights with
the same dimensions as \code{titre}. Column \code{i} is passed as per-observation
weights for biomarker \code{i}, multiplying each observation's log-likelihood
contribution in the Stan model. A plain vector of length \code{nrow(titre)} is
also accepted and broadcast across all biomarkers.
If \code{NULL} (the default) all observations receive equal weight.

\end{description}


\end{SubSubSection}

%
\begin{SubSubSection}{Returns}
A new SeroCOPMulti object
\end{SubSubSection}

\end{SubSection}



\hypertarget{method-SeroCOPMulti-fit_all}{}
%
\begin{SubSection}{Method \code{fit\_all()}}
Fit models for all biomarkers
%
\begin{SubSubSection}{Usage}
\begin{alltt}SeroCOPMulti$fit_all(
  chains = 4,
  iter = 2000,
  warmup = floor(iter/2),
  cores = 1,
  ...
)\end{alltt}

\end{SubSubSection}


%
\begin{SubSubSection}{Arguments}

\begin{description}

\item[\code{chains}] Number of MCMC chains (default: 4)

\item[\code{iter}] Number of iterations per chain (default: 2000)

\item[\code{warmup}] Number of warmup iterations (default: iter/2)

\item[\code{cores}] Number of cores for parallel processing (default: 1)

\item[\code{...}] Additional arguments passed to brms::brm

\end{description}


\end{SubSubSection}

%
\begin{SubSubSection}{Returns}
Self (invisibly)
\end{SubSubSection}

\end{SubSection}



\hypertarget{method-SeroCOPMulti-compare_biomarkers}{}
%
\begin{SubSection}{Method \code{compare\_biomarkers()}}
Compare biomarkers on AUC and LOO metrics
%
\begin{SubSubSection}{Usage}
\begin{alltt}SeroCOPMulti$compare_biomarkers()\end{alltt}

\end{SubSubSection}


%
\begin{SubSubSection}{Returns}
Data frame with comparison metrics
\end{SubSubSection}

\end{SubSection}



\hypertarget{method-SeroCOPMulti-plot_comparison}{}
%
\begin{SubSection}{Method \code{plot\_comparison()}}
Plot biomarkers on AUC vs LOO plane
%
\begin{SubSubSection}{Usage}
\begin{alltt}SeroCOPMulti$plot_comparison(add_labels = TRUE)\end{alltt}

\end{SubSubSection}


%
\begin{SubSubSection}{Arguments}

\begin{description}

\item[\code{add\_labels}] Logical, add biomarker labels (default: TRUE)

\end{description}


\end{SubSubSection}

%
\begin{SubSubSection}{Returns}
A ggplot object
\end{SubSubSection}

\end{SubSection}



\hypertarget{method-SeroCOPMulti-plot_all_curves}{}
%
\begin{SubSection}{Method \code{plot\_all\_curves()}}
Plot all fitted curves
%
\begin{SubSubSection}{Usage}
\begin{alltt}SeroCOPMulti$plot_all_curves()\end{alltt}

\end{SubSubSection}


%
\begin{SubSubSection}{Returns}
A ggplot object
\end{SubSubSection}

\end{SubSection}



\hypertarget{method-SeroCOPMulti-extract_group_parameters}{}
%
\begin{SubSection}{Method \code{extract\_group\_parameters()}}
Extract group-specific parameters for all biomarkers (hierarchical models only)
%
\begin{SubSubSection}{Usage}
\begin{alltt}SeroCOPMulti$extract_group_parameters()\end{alltt}

\end{SubSubSection}


%
\begin{SubSubSection}{Returns}
A data.frame with group-specific parameter estimates for each biomarker
\end{SubSubSection}

\end{SubSection}



\hypertarget{method-SeroCOPMulti-plot_group_curves}{}
%
\begin{SubSection}{Method \code{plot\_group\_curves()}}
Plot group-specific curves for all biomarkers (hierarchical models only)
%
\begin{SubSubSection}{Usage}
\begin{alltt}SeroCOPMulti$plot_group_curves(
  title = "Group-Specific Correlates by Biomarker"
)\end{alltt}

\end{SubSubSection}


%
\begin{SubSubSection}{Arguments}

\begin{description}

\item[\code{title}] Plot title

\end{description}


\end{SubSubSection}

%
\begin{SubSubSection}{Returns}
A ggplot object (returns NULL if not a hierarchical model)
\end{SubSubSection}

\end{SubSection}



\hypertarget{method-SeroCOPMulti-extract_cop_multi}{}
%
\begin{SubSection}{Method \code{extract\_cop\_multi()}}
Extract the correlate of protection conditional on exposure for all biomarkers.
%
\begin{SubSubSection}{Usage}
\begin{alltt}SeroCOPMulti$extract_cop_multi(correlate_of_risk_list, upper_bound = 0.7)\end{alltt}

\end{SubSubSection}


%
\begin{SubSubSection}{Arguments}

\begin{description}

\item[\code{correlate\_of\_risk\_list}] List of numeric vectors of correlates of risk for each biomarker.

\item[\code{upper\_bound}] Numeric value for the upper bound (default: 0.7).

\end{description}


\end{SubSubSection}

%
\begin{SubSubSection}{Returns}
List of numeric vectors of correlates of protection for each biomarker.
\end{SubSubSection}

\end{SubSection}



\hypertarget{method-SeroCOPMulti-clone}{}
%
\begin{SubSection}{Method \code{clone()}}
The objects of this class are cloneable with this method.
%
\begin{SubSubSection}{Usage}
\begin{alltt}SeroCOPMulti$clone(deep = FALSE)\end{alltt}

\end{SubSubSection}


%
\begin{SubSubSection}{Arguments}

\begin{description}

\item[\code{deep}] Whether to make a deep clone.

\end{description}


\end{SubSubSection}

\end{SubSection}

\end{Section}
%
\begin{Examples}
\begin{ExampleCode}
## Not run: 
# Create multi-biomarker data
titres <- matrix(rnorm(600), ncol = 3)
colnames(titres) <- c("IgG", "IgA", "Neutralization")
infected <- rbinom(200, 1, 0.3)

# Initialize and fit
multi_model <- SeroCOPMulti$new(titre = titres, infected = infected)
multi_model$fit_all(chains = 4, iter = 2000)

# Compare biomarkers
multi_model$compare_biomarkers()
multi_model$plot_comparison()

# With hierarchical effects
age_group <- sample(c("Young", "Middle", "Old"), 200, replace = TRUE)
hier_model <- SeroCOPMulti$new(titre = titres, infected = infected, group = age_group)
hier_model$fit_all(chains = 4, iter = 2000)
hier_model$plot_group_curves()
hier_model$extract_group_parameters()

## End(Not run)
\end{ExampleCode}
\end{Examples}
\HeaderA{simulate\_serocop\_data}{Simulate data for correlates of protection analysis}{simulate.Rul.serocop.Rul.data}
%
\begin{Description}
Generate synthetic data from a four-parameter logistic model for testing
and validation purposes.
\end{Description}
%
\begin{Usage}
\begin{verbatim}
simulate_serocop_data(
  n = 200,
  floor = 0.05,
  ceiling = 0.9,
  ec50 = 1.5,
  slope = 2,
  titre_mean = 2,
  titre_sd = 1.5,
  seed = NULL
)
\end{verbatim}
\end{Usage}
%
\begin{Arguments}
\begin{ldescription}
\item[\code{n}] Number of observations

\item[\code{floor}] Lower asymptote (minimum infection probability)

\item[\code{ceiling}] Upper asymptote (maximum infection probability)

\item[\code{ec50}] Titre at 50\% between floor and ceiling

\item[\code{slope}] Steepness of the curve

\item[\code{titre\_mean}] Mean of titre distribution (normal)

\item[\code{titre\_sd}] Standard deviation of titre distribution

\item[\code{seed}] Random seed for reproducibility
\end{ldescription}
\end{Arguments}
%
\begin{Value}
A list containing:
\begin{ldescription}
\item[\code{titre}] Vector of simulated titres
\item[\code{infected}] Vector of simulated infection outcomes
\item[\code{prob\_true}] True infection probabilities
\item[\code{params}] List of true parameters used for simulation
\end{ldescription}
\end{Value}
%
\begin{Examples}
\begin{ExampleCode}
# Simulate data with default parameters
data <- simulate_serocop_data(n = 200, seed = 123)

# Custom parameters
data <- simulate_serocop_data(
  n = 500,
  floor = 0.02,
  ceiling = 0.95,
  ec50 = 2.0,
  slope = 3.0,
  seed = 456
)
\end{ExampleCode}
\end{Examples}
\printindex{}
\end{document}
